% !TeX root = ../main.tex

\chapter*{Resum}
\addcontentsline{toc}{chapter}{Resum}

La síntesi de respostes impulsionals de sala (RIR) completes mitjançant simulació acústica convencional presenta un cost computacional elevat que en limita l'aplicació pràctica en temps real. Aquest treball parteix de la hipòtesi que la porció inicial de 50~ms d'una RIR, on es concentren les reflexions primerenques, conté informació geomètrica suficient per inferir i extrapolar les propietats estadístiques de la seva cua reverberant tardana. Per abordar aquesta tasca, s'ha implementat i adaptat el model codificador-decodificador DECOR, fonamentat en el processament de senyal diferenciable per predir les envolupants de decaïment multiexponencial. El sistema s'ha entrenat i validat  utilitzant dades de sales completament noves i no vistes per la xarxa neuronal, procedents del corpus sintètic BIRD. Els resultats demostren una convergència estable i una alta capacitat de reconstrucció de la dinàmica temporal global, capturant correctament descriptors físics essencials: un error absolut mitjà de la EDC (EDC MAE) de 5.39~dB, una arrel de l'error quadràtic mitjà de la EDC (EDC RMSE) de 7.26~dB, un error percentual absolut mitjà del temps de reverberació ($T_{60}$ MAPE) de 8.1\% i un error quadràtic mitjà de la relació directe-reverberant (DRR MSE) de 7.24~dB. Tanmateix, les anàlisis qualitatives mostren que la fidelitat espectral i temporal es veu compromesa, fet que es tradueix en una pèrdua de la microestructura fina en el rang d'altes freqüències. Es conclou que, malgrat aquestes limitacions espectrals compartides amb el model original davant dades no vistes, la recerca valida la hipòtesi de partida. L'enfocament destaca per la seva eficiència computacional tant respecte als mètodes de simulació geomètrica convencional, com respecte a altres xarxes neuronals destinades a la mateixa tasca, i consolida aquesta arquitectura com una eina eficaç per al prototipat acústic ràpid i l'augment de dades.

\vfill % <--- Añade esto al final del texto
\chapter*{Resumen}
\addcontentsline{toc}{chapter}{Resumen}

La síntesis de respuestas impulsionales de sala (RIR) completas mediante simulación acústica convencional presenta un coste computacional elevado que limita su aplicación práctica en tiempo real. Este trabajo parte de la hipótesis de que el segmento inicial de 50~ms de una RIR, donde se concentran las  reflexiones tempranas, contiene información geométrica suficiente para inferir y extrapolar las propiedades estadísticas de su cola reverberante tardía. Para abordar esta tarea, se ha implementado y adaptado el modelo codificador-decodificador DECOR, el cual se fundamenta en el procesamiento de señal diferenciable para predecir las envolventes de decaimiento multi-exponencial. El sistema ha sido entrenado y validado utilizando datos de salas completamente nuevas y no vistas por la red neuronal, procedentes del corpus sintético BIRD. Los resultados demuestran una convergencia estable y una alta capacidad de reconstrucción de la dinámica temporal global, capturando correctamente descriptores físicos esenciales: un error absoluto medio de la EDC (EDC MAE) de 5.39~dB, una raíz del error cuadrático medio de la EDC (EDC RMSE) de 7.26~dB, un error porcentual absoluto medio del tiempo de reverberación ($T_{60}$ MAPE) del 8.1\% y un error cuadrático medio de la relación directo-reverberante (DRR MSE) de 7.24~dB. No obstante, los análisis cualitativos muestran que la fidelidad espectral y temporal se ve comprometida, lo que se traduce en una pérdida de la microestructura fina en el rango de altas frecuencias. Se concluye que, a pesar de estas limitaciones espectrales compartidas con el modelo original ante datos no vistos, la investigación valida la hipótesis de partida. El enfoque destaca por su eficiencia computacional tanto frente a los métodos de simulación geométrica convencional como frente a otras redes neuronales destinadas a la misma tarea, consolidando esta arquitectura como una herramienta eficaz para el prototipado acústico rápido y el aumento de datos.
\vfill % <--- Añade esto al final del texto
\chapter*{Abstract}
\addcontentsline{toc}{chapter}{Abstract}

The synthesis of complete room impulse responses (RIRs) through conventional acoustic simulation carries a high computational cost that limits practical real-time deployment. This study is built upon the hypothesis that the initial 50 ms segment of an RIR, which encompasses the early reflections, holds sufficient geometric information to statistically infer and extrapolate the properties of its late reverberation tail. To address this task, we implemented and adapted the DECOR encoder-decoder model, which is grounded in differentiable signal processing to predict multi-exponential decay envelopes. The system was trained and validated using data from completely new rooms unseen by the neural network, drawn from the synthetic BIRD corpus. Results show stable convergence and high capability of reconstructing the global temporal dynamics, correctly capturing essential physical descriptors: an EDC mean absolute error (EDC MAE) of 5.39~dB, an EDC root mean square error (EDC RMSE) of 7.26~dB, a mean absolute percentage error for reverberation time ($T_{60}$ MAPE) of 8.1\%, and a direct-to-reverberant ratio mean square error (DRR MSE) of 7.24~dB. However, qualitative analyses show that spectral and temporal fidelity are compromised, resulting in a loss of fine microstructure in the high-frequency range. We conclude that, despite these spectral limitations, which are also shared by the original model on unseen data, the study validates the initial hypothesis. The approach stands out for its computational efficiency compared to both conventional geometric simulation methods and other neural networks designed for the same task, consolidating this architecture as an effective tool for rapid acoustic prototyping and data augmentation.
\vfill % <--- Añade esto al final del texto
